\section{Limitaciones del Proyecto}

Todo modelo predictivo tiene un techo definido por los datos y los recursos con que
se construye. Reconocer esas restricciones desde el inicio no es una debilidad del
proyecto: es una condición para interpretar correctamente sus resultados.

\subsection{Limitaciones Técnicas}

El entrenamiento se hará en Google Colab (GPU compartida) y equipos personales.
Eso impone dos restricciones reales que vale la pena nombrar sin rodeos.

Primero, la memoria disponible por sesión puede ser insuficiente para procesar el
total de registros de SECOP~II de una sola vez. Según el volumen final de datos,
puede ser necesario recurrir a submuestreo estratificado o dividir el procesamiento
en lotes. Ninguna de las dos opciones invalida el experimento, pero sí agrega
un paso de decisión metodológica que habrá que documentar.

Segundo, los espacios de búsqueda de hiperparámetros tendrán que acotarse.
Una búsqueda exhaustiva sobre XGBoost con muchos parámetros puede tomar horas
en los entornos disponibles, lo que hace inviable esa estrategia. La búsqueda
aleatoria (\texttt{RandomizedSearchCV}) es la alternativa estándar en estos casos
y produce resultados competitivos con una fracción del costo computacional, aunque
sin garantía de encontrar la configuración óptima global.

Para que los experimentos sean reproducibles, se fijarán semillas aleatorias y se
registrarán las versiones exactas de todas las librerías en un archivo de entorno.

\subsection{Limitaciones de los Datos}

Las restricciones de datos son más difíciles de mitigar que las técnicas, porque
dependen de cómo está construida la plataforma SECOP~II y de qué información
el Estado colombiano publica de forma abierta.

\subsubsection{No existe una etiqueta oficial de incumplimiento}

SECOP~II no tiene un campo que diga si un contrato terminó mal. Para construir
la variable objetivo hay que inferirla: un contrato se etiqueta como problemático
si acumula varias adiciones, muestra suspensiones reiteradas o si la diferencia
entre la fecha de terminación pactada y la fecha real de liquidación supera un
umbral razonable. Esta aproximación introduce ruido en las etiquetas. Algunos
contratos pueden quedar mal clasificados, lo que afecta la calidad del
entrenamiento supervisado y hace que las métricas del modelo sean probablemente
algo optimistas respecto a un escenario con etiquetas perfectas.

\subsubsection{Calidad irregular de los registros}

Los datos de SECOP~II no son uniformes. Hay campos vacíos, registros duplicados,
municipios codificados de maneras distintas según la entidad que los ingresó y
fechas con formatos inconsistentes. El pipeline de limpieza se encargará de la
mayoría de estos problemas, pero es probable que un porcentaje de contratos deba
descartarse por falta de información mínima. Si ese porcentaje no es aleatorio
---por ejemplo, si las entidades con peor calidad de datos son también las más
propensas a irregularidades--- el conjunto de entrenamiento quedará sesgado.

\subsubsection{El modelo no ve lo que no está publicado}

Al trabajar exclusivamente con datos abiertos, el modelo queda ciego ante
irregularidades que no dejan rastro en SECOP~II. Vínculos societarios entre
contratistas, historial de sanciones en otras entidades o transferencias fuera
del sistema no son capturables con este enfoque. Eso no hace inútil al modelo:
hay suficiente evidencia internacional de que las variables contractuales públicas
sí predicen riesgo. Pero sí fija un límite claro a lo que puede detectar.